% ---------------------------------------------------------------------------
% Artículo científico de aplicación - Finanzas I (Unidad II)
% Daniela Andrea Vásquez Medina - 2024200537D - Tema 50
% Sección 2: Materiales y métodos. Apartados 2.1 y 2.2 (borrador 2).
% Numeración tomada del artículo de la Unidad I: Antecedentes = 1, Métodos = 2.
%
% Se incluye en el documento principal con % ---------------------------------------------------------------------------
% Artículo científico de aplicación - Finanzas I (Unidad II)
% Daniela Andrea Vásquez Medina - 2024200537D - Tema 50
% Sección 2: Materiales y métodos. Apartados 2.1 y 2.2 (borrador 2).
% Numeración tomada del artículo de la Unidad I: Antecedentes = 1, Métodos = 2.
%
% Se incluye en el documento principal con % ---------------------------------------------------------------------------
% Artículo científico de aplicación - Finanzas I (Unidad II)
% Daniela Andrea Vásquez Medina - 2024200537D - Tema 50
% Sección 2: Materiales y métodos. Apartados 2.1 y 2.2 (borrador 2).
% Numeración tomada del artículo de la Unidad I: Antecedentes = 1, Métodos = 2.
%
% Se incluye en el documento principal con % ---------------------------------------------------------------------------
% Artículo científico de aplicación - Finanzas I (Unidad II)
% Daniela Andrea Vásquez Medina - 2024200537D - Tema 50
% Sección 2: Materiales y métodos. Apartados 2.1 y 2.2 (borrador 2).
% Numeración tomada del artículo de la Unidad I: Antecedentes = 1, Métodos = 2.
%
% Se incluye en el documento principal con \input{secciones/03_metodos}
% Las citas usan las claves de articulo/referencias.bib
% ---------------------------------------------------------------------------

\section{Materiales y métodos}

\subsection{Enfoque, tipo, nivel y diseño}

La investigación es de tipo aplicada, porque no busca generar teoría nueva sino
emplear herramientas de evaluación financiera ya establecidas para resolver un
problema concreto de decisión. El enfoque es cuantitativo, ya que todas las
conclusiones se derivan del cálculo de indicadores numéricos sobre datos
observados. El diseño es no experimental, dado que las variables no se
manipulan: el precio del cobre y el tipo de cambio se toman tal como ocurrieron
en el mercado. Por último, el alcance temporal es longitudinal, porque se trabaja
con una serie de precios diarios a lo largo de dieciséis años, entre 2010 y 2025.

La unidad de análisis es un proyecto minero de cobre a tajo abierto, construido
con parámetros físicos y financieros tomados de fuentes públicas verificables.
Conviene precisar por qué se modela un proyecto de este tipo y no una mina en
operación. Las empresas en producción no publican sus flujos de caja por unidad
minera con el detalle que exige un cálculo de VAN y TIR, de modo que trabajar con
una de ellas obligaría a estimar la mayor parte de los datos. En cambio, los
estudios económicos preliminares de proyectos aún no desarrollados sí publican
abiertamente la escala, la ley, la recuperación y la inversión estimada. Modelar
un proyecto tipo a partir de esa información permite que cada supuesto tenga una
fuente declarada y que cualquier lector pueda rehacer el cálculo.

Esta decisión tiene además respaldo en la literatura. Trabajar con un proyecto
hipotético representativo es una estrategia aceptada en la economía minera:
\textcite{cunningham2025} evalúa de manera cualitativa los perfiles de riesgo y
retorno del financiamiento de proyectos comparando un depósito de pórfido de cobre
tipo con un depósito kimberlítico diamantífero tipo, y analiza cuatro etapas
sucesivas: el depósito, la construcción, la producción y la venta final del mineral.
El proyecto modelado en este trabajo corresponde precisamente a un pórfido de cobre.

Esa literatura también advierte por qué conviene tratar con cautela las cifras de los
estudios preliminares. \textcite{lagos2026} evaluaron el cumplimiento de la propuesta
de valor de 80 proyectos de cobre nuevos desarrollados entre 2015 y 2022, y
encontraron que, en promedio, presentaron un sobrecosto de 21,1\,\% y un retraso de
3,3 años respecto de las estimaciones de prefactibilidad. Entre las causas recurrentes
señalan conflictos comunitarios no resueltos, dificultades de permisos ambientales,
infraestructura de soporte insuficiente, inflación de costos e interrupciones de
financiamiento. Este antecedente es relevante para el presente estudio porque sus
parámetros provienen de un estudio económico preliminar, es decir, de una etapa aún
más temprana que la prefactibilidad.

\subsection{El proyecto modelado}

Los parámetros físicos provienen del estudio económico preliminar actualizado del
proyecto Cotabambas, ubicado en la región de Apurímac, difundido por Panoro
Minerals en su comunicado oficial del 22 de septiembre de 2015. Se adoptó el
criterio de tomar la escala, la vida útil y la ley de una sola fuente, en lugar de
combinar cifras de documentos distintos. La razón es de coherencia interna: la
capacidad de procesamiento, los años de operación y la ley de alimentación de un
proyecto se calculan conjuntamente en el plan minero, de modo que mezclarlos con
los de otro estudio produciría un proyecto que no existe. Este criterio se aplicó
incluso cuando existía información más reciente, como el informe técnico NI 43-101
de Cotabambas de 2023, cuyos recursos actualizados no se usaron para el caso base
justamente porque no vienen acompañados del plan minero correspondiente.

La capacidad de procesamiento es de 80\,000 toneladas por día y la vida útil del
proyecto es de 17 años, lo que equivale a un total de 483 millones de toneladas de
mineral procesado. En el modelo se distribuye ese tonelaje en partes iguales entre
los diecisiete años, de modo que la capacidad permanece constante y la vida útil
varía solo cuando se la somete a sensibilidad.

La ley de alimentación de referencia es de 0,318\,\% de cobre. No se tomó de una
tabla directamente, sino que se calculó como el promedio ponderado por tonelaje de
los dos recursos que reporta el propio estudio preliminar: 127,3 millones de
toneladas indicadas con 0,37\,\% de cobre y 355,8 millones de toneladas inferidas
con 0,30\,\%. El cálculo es el siguiente:

\begin{equation}
\text{Ley}_{\text{ponderada}} =
\frac{127{,}3 \times 0{,}37 + 355{,}8 \times 0{,}30}{483{,}1} \approx 0{,}318\,\%
\end{equation}

Se optó por la ley ponderada y no por la del recurso indicado por sí solo porque
la vida útil de diecisiete años solo se alcanza procesando ambos recursos, de
manera que la ley que efectivamente alimentaría la planta es la del conjunto.

La recuperación metalúrgica adoptada es de 84\,\%. Este parámetro constituye la
principal transferencia desde otra fuente y así se declara: corresponde a la
recuperación de cobre del pórfido fresco reportada en la Tabla 13-1 del informe
técnico NI 43-101 del proyecto La Granja, de First Quantum Minerals, publicado en
mayo de 2026. Se eligió ese valor porque el pórfido fresco es mineral sulfurado
primario, equivalente al material hipógeno que domina el recurso de Cotabambas, y
porque el mismo informe señala que la mineralización pórfida sulfurada primaria
alcanza recuperaciones de entre 80\,\% y 85\,\%. La limitación de usar un dato de
otro proyecto se declara expresamente en el apartado de limitaciones.

Con esos tres parámetros, la producción anual de cobre recuperado del modelo
asciende a 167,3 millones de libras, resultado de multiplicar el tonelaje anual
por la ley y por la recuperación, y de convertir toneladas métricas a libras con
el factor de 2\,204,62 libras por tonelada.

Conviene precisar qué mide esa cifra. Las 167,3 millones de libras corresponden a
\textbf{cobre recuperado}, es decir, el metal que llega al concentrado. No equivalen al
cobre que la fundición efectivamente paga. El estudio económico preliminar reporta
155,1 millones de libras anuales de \textbf{cobre pagable}, una cifra alrededor de
8\,\% menor, y esa diferencia se descompone en dos causas independientes.

La primera es la tasa de pago del concentrado. El informe técnico declara un
\emph{pay factor} para el cobre de entre 96,3\,\% y 96,5\,\%, según se tome la tabla de
términos de fundición o el cuadro de supuestos económicos, pues el documento consigna
ambos valores. Aplicando esa tasa, las 167,3 millones de libras recuperadas del modelo
equivalen a unos 161,1 millones de libras pagables, lo que explica alrededor de
3,7 puntos porcentuales de la diferencia.

La segunda es la recuperación metalúrgica. Despejando la recuperación que igualaría la
producción del modelo con la del estudio preliminar, manteniendo el tonelaje, la ley y
la tasa de pago, se obtiene un valor implícito cercano a 80,8\,\%, por debajo del 84\,\%
adoptado. Este valor implícito es un cálculo propio y no una cifra publicada por el
estudio, que no declara su recuperación en el documento consultado. La comprobación
cierra: con una recuperación de 80,8\,\% y una tasa de pago de 96,3\,\%, el modelo
produce 155,0 millones de libras pagables, frente a las 155,1 reportadas.

Se optó por mantener el 84\,\% porque proviene de una fuente publicada y verificable,
mientras que el 80,8\,\% es una inferencia propia. La consecuencia se declara
expresamente: al trabajar con cobre recuperado y con una recuperación algo superior a la
implícita en el estudio de referencia, los ingresos del modelo constituyen una
\textbf{cota superior}, y el VAN obtenido debe leerse en ese sentido.

% Las citas usan las claves de articulo/referencias.bib
% ---------------------------------------------------------------------------

\section{Materiales y métodos}

\subsection{Enfoque, tipo, nivel y diseño}

La investigación es de tipo aplicada, porque no busca generar teoría nueva sino
emplear herramientas de evaluación financiera ya establecidas para resolver un
problema concreto de decisión. El enfoque es cuantitativo, ya que todas las
conclusiones se derivan del cálculo de indicadores numéricos sobre datos
observados. El diseño es no experimental, dado que las variables no se
manipulan: el precio del cobre y el tipo de cambio se toman tal como ocurrieron
en el mercado. Por último, el alcance temporal es longitudinal, porque se trabaja
con una serie de precios diarios a lo largo de dieciséis años, entre 2010 y 2025.

La unidad de análisis es un proyecto minero de cobre a tajo abierto, construido
con parámetros físicos y financieros tomados de fuentes públicas verificables.
Conviene precisar por qué se modela un proyecto de este tipo y no una mina en
operación. Las empresas en producción no publican sus flujos de caja por unidad
minera con el detalle que exige un cálculo de VAN y TIR, de modo que trabajar con
una de ellas obligaría a estimar la mayor parte de los datos. En cambio, los
estudios económicos preliminares de proyectos aún no desarrollados sí publican
abiertamente la escala, la ley, la recuperación y la inversión estimada. Modelar
un proyecto tipo a partir de esa información permite que cada supuesto tenga una
fuente declarada y que cualquier lector pueda rehacer el cálculo.

Esta decisión tiene además respaldo en la literatura. Trabajar con un proyecto
hipotético representativo es una estrategia aceptada en la economía minera:
\textcite{cunningham2025} evalúa de manera cualitativa los perfiles de riesgo y
retorno del financiamiento de proyectos comparando un depósito de pórfido de cobre
tipo con un depósito kimberlítico diamantífero tipo, y analiza cuatro etapas
sucesivas: el depósito, la construcción, la producción y la venta final del mineral.
El proyecto modelado en este trabajo corresponde precisamente a un pórfido de cobre.

Esa literatura también advierte por qué conviene tratar con cautela las cifras de los
estudios preliminares. \textcite{lagos2026} evaluaron el cumplimiento de la propuesta
de valor de 80 proyectos de cobre nuevos desarrollados entre 2015 y 2022, y
encontraron que, en promedio, presentaron un sobrecosto de 21,1\,\% y un retraso de
3,3 años respecto de las estimaciones de prefactibilidad. Entre las causas recurrentes
señalan conflictos comunitarios no resueltos, dificultades de permisos ambientales,
infraestructura de soporte insuficiente, inflación de costos e interrupciones de
financiamiento. Este antecedente es relevante para el presente estudio porque sus
parámetros provienen de un estudio económico preliminar, es decir, de una etapa aún
más temprana que la prefactibilidad.

\subsection{El proyecto modelado}

Los parámetros físicos provienen del estudio económico preliminar actualizado del
proyecto Cotabambas, ubicado en la región de Apurímac, difundido por Panoro
Minerals en su comunicado oficial del 22 de septiembre de 2015. Se adoptó el
criterio de tomar la escala, la vida útil y la ley de una sola fuente, en lugar de
combinar cifras de documentos distintos. La razón es de coherencia interna: la
capacidad de procesamiento, los años de operación y la ley de alimentación de un
proyecto se calculan conjuntamente en el plan minero, de modo que mezclarlos con
los de otro estudio produciría un proyecto que no existe. Este criterio se aplicó
incluso cuando existía información más reciente, como el informe técnico NI 43-101
de Cotabambas de 2023, cuyos recursos actualizados no se usaron para el caso base
justamente porque no vienen acompañados del plan minero correspondiente.

La capacidad de procesamiento es de 80\,000 toneladas por día y la vida útil del
proyecto es de 17 años, lo que equivale a un total de 483 millones de toneladas de
mineral procesado. En el modelo se distribuye ese tonelaje en partes iguales entre
los diecisiete años, de modo que la capacidad permanece constante y la vida útil
varía solo cuando se la somete a sensibilidad.

La ley de alimentación de referencia es de 0,318\,\% de cobre. No se tomó de una
tabla directamente, sino que se calculó como el promedio ponderado por tonelaje de
los dos recursos que reporta el propio estudio preliminar: 127,3 millones de
toneladas indicadas con 0,37\,\% de cobre y 355,8 millones de toneladas inferidas
con 0,30\,\%. El cálculo es el siguiente:

\begin{equation}
\text{Ley}_{\text{ponderada}} =
\frac{127{,}3 \times 0{,}37 + 355{,}8 \times 0{,}30}{483{,}1} \approx 0{,}318\,\%
\end{equation}

Se optó por la ley ponderada y no por la del recurso indicado por sí solo porque
la vida útil de diecisiete años solo se alcanza procesando ambos recursos, de
manera que la ley que efectivamente alimentaría la planta es la del conjunto.

La recuperación metalúrgica adoptada es de 84\,\%. Este parámetro constituye la
principal transferencia desde otra fuente y así se declara: corresponde a la
recuperación de cobre del pórfido fresco reportada en la Tabla 13-1 del informe
técnico NI 43-101 del proyecto La Granja, de First Quantum Minerals, publicado en
mayo de 2026. Se eligió ese valor porque el pórfido fresco es mineral sulfurado
primario, equivalente al material hipógeno que domina el recurso de Cotabambas, y
porque el mismo informe señala que la mineralización pórfida sulfurada primaria
alcanza recuperaciones de entre 80\,\% y 85\,\%. La limitación de usar un dato de
otro proyecto se declara expresamente en el apartado de limitaciones.

Con esos tres parámetros, la producción anual de cobre recuperado del modelo
asciende a 167,3 millones de libras, resultado de multiplicar el tonelaje anual
por la ley y por la recuperación, y de convertir toneladas métricas a libras con
el factor de 2\,204,62 libras por tonelada.

Conviene precisar qué mide esa cifra. Las 167,3 millones de libras corresponden a
\textbf{cobre recuperado}, es decir, el metal que llega al concentrado. No equivalen al
cobre que la fundición efectivamente paga. El estudio económico preliminar reporta
155,1 millones de libras anuales de \textbf{cobre pagable}, una cifra alrededor de
8\,\% menor, y esa diferencia se descompone en dos causas independientes.

La primera es la tasa de pago del concentrado. El informe técnico declara un
\emph{pay factor} para el cobre de entre 96,3\,\% y 96,5\,\%, según se tome la tabla de
términos de fundición o el cuadro de supuestos económicos, pues el documento consigna
ambos valores. Aplicando esa tasa, las 167,3 millones de libras recuperadas del modelo
equivalen a unos 161,1 millones de libras pagables, lo que explica alrededor de
3,7 puntos porcentuales de la diferencia.

La segunda es la recuperación metalúrgica. Despejando la recuperación que igualaría la
producción del modelo con la del estudio preliminar, manteniendo el tonelaje, la ley y
la tasa de pago, se obtiene un valor implícito cercano a 80,8\,\%, por debajo del 84\,\%
adoptado. Este valor implícito es un cálculo propio y no una cifra publicada por el
estudio, que no declara su recuperación en el documento consultado. La comprobación
cierra: con una recuperación de 80,8\,\% y una tasa de pago de 96,3\,\%, el modelo
produce 155,0 millones de libras pagables, frente a las 155,1 reportadas.

Se optó por mantener el 84\,\% porque proviene de una fuente publicada y verificable,
mientras que el 80,8\,\% es una inferencia propia. La consecuencia se declara
expresamente: al trabajar con cobre recuperado y con una recuperación algo superior a la
implícita en el estudio de referencia, los ingresos del modelo constituyen una
\textbf{cota superior}, y el VAN obtenido debe leerse en ese sentido.

% Las citas usan las claves de articulo/referencias.bib
% ---------------------------------------------------------------------------

\section{Materiales y métodos}

\subsection{Enfoque, tipo, nivel y diseño}

La investigación es de tipo aplicada, porque no busca generar teoría nueva sino
emplear herramientas de evaluación financiera ya establecidas para resolver un
problema concreto de decisión. El enfoque es cuantitativo, ya que todas las
conclusiones se derivan del cálculo de indicadores numéricos sobre datos
observados. El diseño es no experimental, dado que las variables no se
manipulan: el precio del cobre y el tipo de cambio se toman tal como ocurrieron
en el mercado. Por último, el alcance temporal es longitudinal, porque se trabaja
con una serie de precios diarios a lo largo de dieciséis años, entre 2010 y 2025.

La unidad de análisis es un proyecto minero de cobre a tajo abierto, construido
con parámetros físicos y financieros tomados de fuentes públicas verificables.
Conviene precisar por qué se modela un proyecto de este tipo y no una mina en
operación. Las empresas en producción no publican sus flujos de caja por unidad
minera con el detalle que exige un cálculo de VAN y TIR, de modo que trabajar con
una de ellas obligaría a estimar la mayor parte de los datos. En cambio, los
estudios económicos preliminares de proyectos aún no desarrollados sí publican
abiertamente la escala, la ley, la recuperación y la inversión estimada. Modelar
un proyecto tipo a partir de esa información permite que cada supuesto tenga una
fuente declarada y que cualquier lector pueda rehacer el cálculo.

Esta decisión tiene además respaldo en la literatura. Trabajar con un proyecto
hipotético representativo es una estrategia aceptada en la economía minera:
\textcite{cunningham2025} evalúa de manera cualitativa los perfiles de riesgo y
retorno del financiamiento de proyectos comparando un depósito de pórfido de cobre
tipo con un depósito kimberlítico diamantífero tipo, y analiza cuatro etapas
sucesivas: el depósito, la construcción, la producción y la venta final del mineral.
El proyecto modelado en este trabajo corresponde precisamente a un pórfido de cobre.

Esa literatura también advierte por qué conviene tratar con cautela las cifras de los
estudios preliminares. \textcite{lagos2026} evaluaron el cumplimiento de la propuesta
de valor de 80 proyectos de cobre nuevos desarrollados entre 2015 y 2022, y
encontraron que, en promedio, presentaron un sobrecosto de 21,1\,\% y un retraso de
3,3 años respecto de las estimaciones de prefactibilidad. Entre las causas recurrentes
señalan conflictos comunitarios no resueltos, dificultades de permisos ambientales,
infraestructura de soporte insuficiente, inflación de costos e interrupciones de
financiamiento. Este antecedente es relevante para el presente estudio porque sus
parámetros provienen de un estudio económico preliminar, es decir, de una etapa aún
más temprana que la prefactibilidad.

\subsection{El proyecto modelado}

Los parámetros físicos provienen del estudio económico preliminar actualizado del
proyecto Cotabambas, ubicado en la región de Apurímac, difundido por Panoro
Minerals en su comunicado oficial del 22 de septiembre de 2015. Se adoptó el
criterio de tomar la escala, la vida útil y la ley de una sola fuente, en lugar de
combinar cifras de documentos distintos. La razón es de coherencia interna: la
capacidad de procesamiento, los años de operación y la ley de alimentación de un
proyecto se calculan conjuntamente en el plan minero, de modo que mezclarlos con
los de otro estudio produciría un proyecto que no existe. Este criterio se aplicó
incluso cuando existía información más reciente, como el informe técnico NI 43-101
de Cotabambas de 2023, cuyos recursos actualizados no se usaron para el caso base
justamente porque no vienen acompañados del plan minero correspondiente.

La capacidad de procesamiento es de 80\,000 toneladas por día y la vida útil del
proyecto es de 17 años, lo que equivale a un total de 483 millones de toneladas de
mineral procesado. En el modelo se distribuye ese tonelaje en partes iguales entre
los diecisiete años, de modo que la capacidad permanece constante y la vida útil
varía solo cuando se la somete a sensibilidad.

La ley de alimentación de referencia es de 0,318\,\% de cobre. No se tomó de una
tabla directamente, sino que se calculó como el promedio ponderado por tonelaje de
los dos recursos que reporta el propio estudio preliminar: 127,3 millones de
toneladas indicadas con 0,37\,\% de cobre y 355,8 millones de toneladas inferidas
con 0,30\,\%. El cálculo es el siguiente:

\begin{equation}
\text{Ley}_{\text{ponderada}} =
\frac{127{,}3 \times 0{,}37 + 355{,}8 \times 0{,}30}{483{,}1} \approx 0{,}318\,\%
\end{equation}

Se optó por la ley ponderada y no por la del recurso indicado por sí solo porque
la vida útil de diecisiete años solo se alcanza procesando ambos recursos, de
manera que la ley que efectivamente alimentaría la planta es la del conjunto.

La recuperación metalúrgica adoptada es de 84\,\%. Este parámetro constituye la
principal transferencia desde otra fuente y así se declara: corresponde a la
recuperación de cobre del pórfido fresco reportada en la Tabla 13-1 del informe
técnico NI 43-101 del proyecto La Granja, de First Quantum Minerals, publicado en
mayo de 2026. Se eligió ese valor porque el pórfido fresco es mineral sulfurado
primario, equivalente al material hipógeno que domina el recurso de Cotabambas, y
porque el mismo informe señala que la mineralización pórfida sulfurada primaria
alcanza recuperaciones de entre 80\,\% y 85\,\%. La limitación de usar un dato de
otro proyecto se declara expresamente en el apartado de limitaciones.

Con esos tres parámetros, la producción anual de cobre recuperado del modelo
asciende a 167,3 millones de libras, resultado de multiplicar el tonelaje anual
por la ley y por la recuperación, y de convertir toneladas métricas a libras con
el factor de 2\,204,62 libras por tonelada.

Conviene precisar qué mide esa cifra. Las 167,3 millones de libras corresponden a
\textbf{cobre recuperado}, es decir, el metal que llega al concentrado. No equivalen al
cobre que la fundición efectivamente paga. El estudio económico preliminar reporta
155,1 millones de libras anuales de \textbf{cobre pagable}, una cifra alrededor de
8\,\% menor, y esa diferencia se descompone en dos causas independientes.

La primera es la tasa de pago del concentrado. El informe técnico declara un
\emph{pay factor} para el cobre de entre 96,3\,\% y 96,5\,\%, según se tome la tabla de
términos de fundición o el cuadro de supuestos económicos, pues el documento consigna
ambos valores. Aplicando esa tasa, las 167,3 millones de libras recuperadas del modelo
equivalen a unos 161,1 millones de libras pagables, lo que explica alrededor de
3,7 puntos porcentuales de la diferencia.

La segunda es la recuperación metalúrgica. Despejando la recuperación que igualaría la
producción del modelo con la del estudio preliminar, manteniendo el tonelaje, la ley y
la tasa de pago, se obtiene un valor implícito cercano a 80,8\,\%, por debajo del 84\,\%
adoptado. Este valor implícito es un cálculo propio y no una cifra publicada por el
estudio, que no declara su recuperación en el documento consultado. La comprobación
cierra: con una recuperación de 80,8\,\% y una tasa de pago de 96,3\,\%, el modelo
produce 155,0 millones de libras pagables, frente a las 155,1 reportadas.

Se optó por mantener el 84\,\% porque proviene de una fuente publicada y verificable,
mientras que el 80,8\,\% es una inferencia propia. La consecuencia se declara
expresamente: al trabajar con cobre recuperado y con una recuperación algo superior a la
implícita en el estudio de referencia, los ingresos del modelo constituyen una
\textbf{cota superior}, y el VAN obtenido debe leerse en ese sentido.

% Las citas usan las claves de articulo/referencias.bib
% ---------------------------------------------------------------------------

\section{Materiales y métodos}

\subsection{Enfoque, tipo, nivel y diseño}

La investigación es de tipo aplicada, porque no busca generar teoría nueva sino
emplear herramientas de evaluación financiera ya establecidas para resolver un
problema concreto de decisión. El enfoque es cuantitativo, ya que todas las
conclusiones se derivan del cálculo de indicadores numéricos sobre datos
observados. El diseño es no experimental, dado que las variables no se
manipulan: el precio del cobre y el tipo de cambio se toman tal como ocurrieron
en el mercado. Por último, el alcance temporal es longitudinal, porque se trabaja
con una serie de precios diarios a lo largo de dieciséis años, entre 2010 y 2025.

El nivel de la investigación es descriptivo. El estudio mide y caracteriza el
comportamiento de los indicadores financieros del proyecto bajo distintos
escenarios de precio y de inversión, pero no contrasta una relación de
dependencia entre variables ni estima parámetros de un modelo de inferencia. Por
esa razón el estudio no formula hipótesis: la pregunta de investigación indaga
\emph{cómo} se comportan el valor actual neto y la tasa interna de retorno frente
a la variabilidad observada del precio, y no constituye una proposición que pueda
ser declarada verdadera o falsa mediante una prueba estadística.

La unidad de análisis es un proyecto minero de cobre a tajo abierto, construido
con parámetros físicos y financieros tomados de fuentes públicas verificables.
Conviene precisar por qué se modela un proyecto de este tipo y no una mina en
operación. Las empresas en producción no publican sus flujos de caja por unidad
minera con el detalle que exige un cálculo de VAN y TIR, de modo que trabajar con
una de ellas obligaría a estimar la mayor parte de los datos. En cambio, los
estudios económicos preliminares de proyectos aún no desarrollados sí publican
abiertamente la escala, la ley, la recuperación y la inversión estimada. Modelar
un proyecto representativo a partir de esa información permite que cada supuesto tenga una
fuente declarada y que cualquier lector pueda rehacer el cálculo.

Esta decisión tiene además respaldo en la literatura. Trabajar con un proyecto
hipotético representativo es una estrategia aceptada en la economía minera:
\textcite{cunningham2025} evalúa de manera cualitativa los perfiles de riesgo y
retorno del financiamiento de proyectos comparando un depósito de pórfido de cobre
tipo con un depósito kimberlítico diamantífero tipo, y analiza cuatro etapas
sucesivas: el depósito, la construcción, la producción y la venta final del mineral.
El proyecto modelado en este trabajo corresponde precisamente a un pórfido de cobre.

Esa literatura también advierte por qué conviene tratar con cautela las cifras de los
estudios preliminares. \textcite{lagos2026} evaluaron el cumplimiento de la propuesta
de valor de 80 proyectos de cobre nuevos desarrollados entre 2015 y 2022, y
encontraron que, en promedio, presentaron un sobrecosto de 21,1\,\% y un retraso de
3,3 años respecto de las estimaciones de prefactibilidad. Entre las causas recurrentes
señalan conflictos comunitarios no resueltos, dificultades de permisos ambientales,
infraestructura de soporte insuficiente, inflación de costos e interrupciones de
financiamiento. Este antecedente es relevante para el presente estudio porque sus
parámetros provienen de un estudio económico preliminar, es decir, de una etapa aún
más temprana que la prefactibilidad.

\subsection{El proyecto modelado}

Los parámetros físicos provienen del estudio económico preliminar actualizado del
proyecto Cotabambas, ubicado en la región de Apurímac, difundido por Panoro
Minerals en su comunicado oficial del 22 de septiembre de 2015. Se adoptó el
criterio de tomar la escala, la vida útil y la ley de una sola fuente, en lugar de
combinar cifras de documentos distintos. La razón es de coherencia interna: la
capacidad de procesamiento, los años de operación y la ley de alimentación de un
proyecto se calculan conjuntamente en el plan minero, de modo que mezclarlos con
los de otro estudio produciría un proyecto que no existe. Este criterio se aplicó
incluso cuando existía información más reciente, como el informe técnico NI 43-101
de Cotabambas de 2023, cuyos recursos actualizados no se usaron para el caso base
justamente porque no vienen acompañados del plan minero correspondiente.

La capacidad de procesamiento es de 80\,000 toneladas por día y la vida útil del
proyecto es de 17 años, lo que equivale a un total de 483 millones de toneladas de
mineral procesado. En el modelo se distribuye ese tonelaje en partes iguales entre
los diecisiete años, de modo que la capacidad permanece constante y la vida útil
varía solo cuando se la somete a sensibilidad.

La ley de alimentación de referencia es de 0,318\,\% de cobre. No se tomó de una
tabla directamente, sino que se calculó como el promedio ponderado por tonelaje de
los dos recursos que reporta el propio estudio preliminar: 127,3 millones de
toneladas indicadas con 0,37\,\% de cobre y 355,8 millones de toneladas inferidas
con 0,30\,\%. El cálculo es el siguiente:

\begin{equation}
\text{Ley}_{\text{ponderada}} =
\frac{127{,}3 \times 0{,}37 + 355{,}8 \times 0{,}30}{483{,}1} \approx 0{,}318\,\%
\end{equation}

Se optó por la ley ponderada y no por la del recurso indicado por sí solo porque
la vida útil de diecisiete años solo se alcanza procesando ambos recursos, de
manera que la ley que efectivamente alimentaría la planta es la del conjunto.

La recuperación metalúrgica adoptada es de 84\,\%. Este parámetro constituye la
principal transferencia desde otra fuente y así se declara: corresponde a la
recuperación de cobre del pórfido fresco reportada en la Tabla 13-1 del informe
técnico NI 43-101 del proyecto La Granja, de First Quantum Minerals, publicado en
mayo de 2026. Se eligió ese valor porque el pórfido fresco es mineral sulfurado
primario, equivalente al material hipógeno que domina el recurso de Cotabambas, y
porque el mismo informe señala que la mineralización pórfida sulfurada primaria
alcanza recuperaciones de entre 80\,\% y 85\,\%. La limitación de usar un dato de
otro proyecto se declara expresamente en el apartado de limitaciones.

Con esos tres parámetros, la producción anual de cobre recuperado del modelo
asciende a 167,3 millones de libras, resultado de multiplicar el tonelaje anual
por la ley y por la recuperación, y de convertir toneladas métricas a libras con
el factor de 2\,204,62 libras por tonelada.

Conviene precisar qué mide esa cifra. Las 167,3 millones de libras corresponden a
\textbf{cobre recuperado}, es decir, el metal que llega al concentrado. No equivalen al
cobre que la fundición efectivamente paga. El estudio económico preliminar reporta
155,1 millones de libras anuales de \textbf{cobre pagable}, una cifra alrededor de
8\,\% menor, y esa diferencia se descompone en dos causas independientes.

La primera es la tasa de pago del concentrado. El informe técnico declara un
\emph{pay factor} para el cobre de entre 96,3\,\% y 96,5\,\%, según se tome la tabla de
términos de fundición o el cuadro de supuestos económicos, pues el documento consigna
ambos valores. Aplicando esa tasa, las 167,3 millones de libras recuperadas del modelo
equivalen a unos 161,1 millones de libras pagables, lo que explica alrededor de
3,7 puntos porcentuales de la diferencia.

La segunda es la recuperación metalúrgica. Despejando la recuperación que igualaría la
producción del modelo con la del estudio preliminar, manteniendo el tonelaje, la ley y
la tasa de pago, se obtiene un valor implícito cercano a 80,8\,\%, por debajo del 84\,\%
adoptado. Este valor implícito es un cálculo propio y no una cifra publicada por el
estudio, que no declara su recuperación en el documento consultado. La comprobación
cierra: con una recuperación de 80,8\,\% y una tasa de pago de 96,3\,\%, el modelo
produce 155,0 millones de libras pagables, frente a las 155,1 reportadas.

Se optó por mantener el 84\,\% porque proviene de una fuente publicada y verificable,
mientras que el 80,8\,\% es una inferencia propia. La consecuencia se declara
expresamente: al trabajar con cobre recuperado y con una recuperación algo superior a la
implícita en el estudio de referencia, los ingresos del modelo constituyen una
\textbf{cota superior}, y el VAN obtenido debe leerse en ese sentido.
